
\section{Méthodes}

\subsection{Mortalité routière}

\subsubsection{Observation de carcasses}

En 2025, des inventaires de carcasses ont été réalisés sur 30 tronçons routiers de 1 km. Chaque tronçon faisait l'objet de deux visites, et durant chaque visite nous procédions à deux inventaires réalisés selon deux méthodes d'échantillonnage distinctes. Les inventaires se déroulaient entre 5 h 00 et 10 h 00, de manière à réduire le biais lié au retrait ou le déplacement des carcasses par les charognards et les véhicules circulant sur la chaussée.

\addM{Comme présenté au volet 1, trois méthodes ont été utilisées pour détecter les carcasses. Les inventaires par caméra embarquée ont été réalisés à l'aide de la même caméra GoPro Hero 11 Black décrite au volet 1 et avec les mêmes paramètres de configuration. La méthode à pied consistait à parcourir l’accotement en marchant. Les observateurs y repéraient, identifiaient et dénombraient toutes les carcasses détectées, y compris celles situées hors de la chaussée, sur l’accotement gravillonné ou dans la végétation herbacée en bordure de route. Sur les routes secondaires, les observateurs parcouraient successivement les deux voies en veillant à ne pas compter deux fois une même carcasse. Dans le cas où les observateurs ne pouvaient pas traverser le tronçon de manière sécuritaire, ils inventoriaient les deux voies à partir d’un seul côté de la route. Les observateurs notaient la date, l’heure, le taxon et les coordonnées géographiques de chaque carcasse à l’aide d’un GPS GARMIN GPSMAP 79s (Garmin, Schaffhausen, Suisse). La troisième méthode de détection des carcasses utilisait le même drone décrit dans le volet 1 avec les mêmes paramètres de configuration et les mêmes contraintes de conditions de vol.}
  

La méthode d'inventaire par caméra embarquée \addM{GoPro} était utilisée à chaque visite. La seconde méthode \addM{utilisée lors d'une visite donnée était soit le} \delM{de chaque visite )} \addM{drone ou le relevé à pied} \addM{et} dépendait des autorisations et des conditions de vol. Par conséquent, bien qu'une visite impliquait deux méthodes distinctes, un tronçon pouvait cumuler \addM{jusqu'à} trois méthodes différentes \delM{au bout de}\addM{après} deux visites (Tableau~\ref{tab:tab_method}).


\begin{longtable}{llrrl}
  \caption{Nombre de visites et méthodes utilisées pour les inventaires de carcasses, par tronçon.}
  \label{tab:tab_method}
  \\
  \toprule
Section & Tronçon & Catégorie & Visites & Methodes \\ 
  \midrule
T01\_E & T01 &   2 &   2 & À pied+GoPro \\ 
  T01\_O & T01 &   2 &   2 & Drone+GoPro \\ 
  T02 & T02 &   0 &   2 & À pied+GoPro \\ 
  T03\_E & T03 &   0 &   2 & À pied+Drone+GoPro \\ 
  T03\_O & T03 &   0 &   2 & GoPro \\ 
  T09\_E & T09 &   0 &   2 & GoPro \\ 
  T09\_O & T09 &   0 &   2 & À pied+GoPro \\ 
  T14\_E & T14 &   0 &   2 & À pied+GoPro \\ 
  T14\_O & T14 &   0 &   2 & À pied+GoPro \\ 
  T15\_E & T15 &   0 &   2 & GoPro \\ 
  T15\_O & T15 &   0 &   2 & À pied+GoPro \\ 
  T20\_E & T20 &   0 &   2 & Drone+GoPro \\ 
  T20\_O & T20 &   0 &   2 & Drone+GoPro \\ 
  T22\_E & T22 &   2 &   2 & À pied+GoPro \\ 
  T22\_O & T22 &   2 &   2 & GoPro \\ 
  T25\_E & T25 &   1 &   2 & À pied+Drone+GoPro \\ 
  T25\_O & T25 &   1 &   2 & GoPro \\ 
  T26\_E & T26 &   0 &   2 & GoPro \\ 
  T26\_O & T26 &   0 &   2 & À pied+GoPro \\ 
  T28 & T28 &   0 &   2 & À pied+Drone+GoPro \\ 
  T30\_N & T30 &   0 &   2 & GoPro \\ 
  T30\_S & T30 &   0 &   2 & À pied+Drone+GoPro \\ 
  T31\_E & T31 &   1 &   2 & À pied+Drone+GoPro \\ 
  T31\_O & T31 &   1 &   2 & GoPro \\ 
  T35 & T35 &   1 &   2 & À pied+Drone+GoPro \\ 
  T36 & T36 &   2 &   2 & À pied+Drone+GoPro \\ 
  T39\_E & T39 &   1 &   2 & GoPro \\ 
  T39\_O & T39 &   1 &   2 & À pied+Drone+GoPro \\ 
  T40\_E & T40 &   2 &   2 & Drone+GoPro \\ 
  T40\_O & T40 &   2 &   2 & GoPro \\ 
  T41\_E & T41 &   2 &   2 & GoPro \\ 
  T41\_O & T41 &   2 &   3 & À pied+Drone+GoPro \\ 
  T42\_E & T42 &   2 &   2 & GoPro \\ 
  T42\_O & T42 &   2 &   3 & À pied+Drone+GoPro \\ 
  T45\_N & T45 &   2 &   2 & À pied+GoPro \\ 
  T45\_S & T45 &   2 &   2 & GoPro \\ 
  T47\_E & T47 &   2 &   2 & Drone+GoPro \\ 
  T47\_O & T47 &   2 &   2 & GoPro \\ 
  T48\_E & T48 &   1 &   2 & À pied+GoPro \\ 
  T48\_O & T48 &   1 &   2 & GoPro \\ 
  T49 & T49 &   2 &   2 & À pied+Drone+GoPro \\ 
  T52\_E & T52 &   2 &   2 & GoPro \\ 
  T52\_O & T52 &   2 &   2 & À pied+GoPro \\ 
  T53\_E & T53 &   1 &   2 & GoPro \\ 
  T53\_O & T53 &   1 &   2 & À pied+Drone+GoPro \\ 
  T57\_E & T57 &   2 &   2 & Drone+GoPro \\ 
  T57\_O & T57 &   2 &   2 & GoPro \\ 
  T59 & T59 &   1 &   2 & À pied+GoPro \\ 
  T61\_N & T61 &   2 &   2 & À pied+Drone+GoPro \\ 
  T61\_S & T61 &   2 &   2 & GoPro \\ 
  T62\_E & T62 &   1 &   2 & GoPro \\ 
  T62\_O & T62 &   1 &   2 & À pied+Drone+GoPro \\ 
  T65\_N & T65 &   1 &   2 & À pied+Drone+GoPro \\ 
  T65\_S & T65 &   1 &   2 & GoPro \\ 
   \bottomrule
\label{tab:visites}
\end{longtable}


\comM{Mieux de garder ça succinct, ça ne sert à rien de s'autoplagier entre différentes sections.}\delM{(GoPro, San Mateo, Californie, États-Unis) fixée au centre de la partie avant du capot du véhicule au moyen d'un trépied magnétique 3-Footed Monster. Afin d'enregistrer la chaussée lors du parcours de chaque tronçon, la caméra était configurée en résolution 4K au format 4:3, à une cadence de 30 images par seconde. La caméra était pilotée depuis l'habitacle du véhicule au moyen de l'application Quik (GoPro, San Mateo, Californie, États-Unis) installée sur un téléphone intelligent. Les carcasses ont été dénombrées a posteriori par visionnement des enregistrements vidéo. La caméra étant équipée d'un récepteur GPS intégré, chaque carcasse détectée à l'écran a pu être géolocalisée en croisant le minutage de sa détection avec les données de position contenues dans les métadonnées des fichiers \texttt{.mp4}. Dans le cas des inventaires par caméra embarquée, les deux voies de circulation étaient systématiquement parcourues. La méthode à pied consistait à parcourir l'accotement en marchant, celui-ci étant constitué de gravel ou enherbé selon le tronçon. Les observateurs y repéraient, identifiaient et dénombraient toutes les carcasses détectées, y compris celles situées hors de la chaussée, sur l'accotement gravillonné ou dans la végétation herbacée en bordure de route. Sur les routes secondaires, les observateurs parcouraient successivement les deux voies en veillant à ne pas compter deux fois une même carcasse. Dans le cas où les observateurs ne pouvaient pas traverser le tronçon de manière sécuritaire, ils inventoriaient les deux voies à partir d'un seul côté de la route. Les observateurs notaient la date, l'heure, le taxon et les coordonnées géographiques de chaque carcasse à l'aide d'un GPS GARMIN GPSMAP 79s (Garmin, Schaffhausen, Suisse). La troisième méthode de détection des carcasses, utilisant un aéronef télépiloté (ici désigné par « drone »), venait compléter les deux méthodes précédemment présentées (suivi à pied avec jumelles et caméra fixée sur véhicule) utilisait le même drone décrit dans le volet 1 avec les mêmes paramètres de configuration et les mêmes contraintes de conditions de vol. Pour rappel, cette méthode pouvait se substituer au suivi à pied dans certaines conditions. Les journées d’échantillonnage avec drone étaient moins fréquentes que celles des deux autres méthodes, car les vols nécessitent des conditions météorologiques particulières : vents inférieurs à 30 km/h et absence de précipitations. Le drone utilisé était un DJI Mavic 3 Enterprise Thermal (DJI, Shenzhen, Chine). Celui-ci a été immatriculé par Transports Canada avant les vols, et sa plaque d'immatriculation est visible sur une des pales de l'engin. Puisque le drone pesait plus de 250 g, il était nécessaire d’obtenir une certification de base pour devenir pilote. Pendant le vol, une à trois autres personnes de l'équipe doivent être des observateurs visuels, afin de conserver un contact visuel avec l’aéronef en tout temps.}
  
\delM{Dès le lever du soleil, le drone était déployé le long de l'A10 et de la R104 une fois toutes les \addM{2--3} semaines lors des conditions optimales de vol. Les vols ont été automatisés pour que le drone effectue des tronçons linéaires, en restant parallèle aux routes et en volant à au moins 5 m de la chaussée. Les trajets de vol étaient de 500 m chacun, le but étant de voler de part et d'autre des futurs passages, soit 250 m avant et 250 m après les points définis de ceux-ci (Figure 3). L’altitude de vol était de 40 m. Nous avons utilisé deux types de caméras, une de grand-angle et une thermique. La caméra grand-angle possède une qualité d’image de 20 mégapixels, et une lentille avec un champ de vision de 84°, un équivalent format de 24 mm, une ouverture de f/2.8 et une mise au point à 1 m. La qualité du 4K permet d’obtenir 60 images par seconde. La caméra thermique de ce drone fonctionne avec un imageur thermique, le Microbolomètre VOx non refroidi et d’un spotmètre. La lentille a un champ de vision de 61°, un équivalent format de 40 mm, une ouverture de f/1.0 et une mise au point à 5 m. La plage de mesure va de -20° à 150°C, avec une précision de +/- 2°C.
Le drone circulait à une vitesse de 2 m/s, afin de bien couvrir la chaussée lors du visionnement. La caméra du drone était orientée à 90 degrés pour faire face à la route et agrandissait l’image continuellement à un grossissement de X2. Les fichiers vidéo des vols de 500 m de qualité 4K (ratio d’image en 4:3) ont été visionnés en laboratoire et les données d'observation sur la route ont été saisies en notant le moment (minutes, secondes) de la séquence vidéo. Les carcasses ont été géolocalisées, au même titre que les deux précédentes méthodes. 
Le but de cette technique était de mettre en relation les données d'inventaire sur la faune retrouvée proche des futurs passages fauniques et les individus retrouvés morts sur la route. }



\subsubsection{Persistance de carcasses sur la route}

Afin d'estimer la durée de présence d'une carcasse avant son enlèvement par les charognards, nous avons réalisé des tests de persistance sur la majorité des tronçons suivis \addM{en suivant le même protocole que décrit dans le volet 1}. \delM{Ainsi, à}\addM{À partir} du mois de juin, deux tests hebdomadaires étaient effectués, chacun sur un tronçon différent. L'un débutait après le lever du soleil (vers 5h00), l'autre un peu avant le coucher du soleil (vers 20h00). L'emplacement (la voie et le kilométrage) de chaque test était tiré aléatoirement. Lors de chaque test, des morceaux de poulet cru de 1.5 cm x 3 cm x 3 cm étaient déposés sur l'accotement et suivis par un piège photographique paramétré pour capturer une image toutes les 15 secondes. La caméra, une Spypoint Force-20, était montée sur un pivot à une hauteur standard de 30 cm du sol, fixé à un piquet de bois lui-même stabilisé par des planches à sa base. Après environ 24 heures suivant son déploiement, la caméra était retirée. L'heure et la date de pose de la caméra, l'heure et la date de retrait, les coordonnées, et le cas échéant le temps écoulé entre la pose et la prédation ainsi que l'identité du charognard, étaient consignés.



\subsubsection{Observations complémentaires}
\comM[inline]{Est-ce pertinent de maintenir une sous-sous-section \og Observations complémentaires \fg{}? Il n'y a rien dans la section pour le moment.}

\subsection{Suivi de la petite faune}

\subsubsection{Pièges photographiques}

Des pièges photographiques (GardePro T5CF, Hong Kong, Chine) ont été installés en mai 2025 à chacun des 30 tronçons sélectionnés. Un appareil a été disposé à proximité d'un ponceau dans chaque tronçon. \addM{Nous avons utilisé les mêmes paramètres de configuration que décrits dans le volet 1.} Ces pièges photographiques ont permis de détecter la faune, et plus particulièrement les mammifères de petite et moyenne taille, ainsi que l'herpétofaune. Ces pièges photographiques ont été retirés en octobre 2025. \delM{Les pièges photographiques ont été programmés avec deux modes de déclenchement: la prise de photo à intervalle régulier de 5 minutes (\emph{time lapse}) pour les ectothermes et la prise d'une photo après détection de mouvement pour les endothermes. Afin de ne pas saturer le stockage des cartes mémoires, nous avons réglé la qualité à 4 mégapixels. Le principal avantage du modèle de caméra que nous avons choisi est la mise au point rapprochée des animaux détectés avec un focus à 20 cm, facilitant l'identification des petits organismes.}

\addM{Nous avons utilisé l'outil MegaDetector pour traiter les images récoltées, en suivant la même approche que pour le volet 1.} \delM{Cet outil permet de trier les images qui contiennent des photos avec animaux et de celles qui n'en contiennent pas. MegaDetector n'identifie pas à l'espèce. C'est pour cette raison que les} \addM{Les images avec détection présumée et avec une confiance $\geq 75$\% ont ensuite été validées par un humain pour déterminer si la détection était avérée et le cas échéant, d'identifier les individus dans les images au taxon le plus bas possible.}
  
\delM{Dans un premier temps, nous avons utilisé l'outil de reconnaissance automatique MegaDetector afin de faire un tri initial des photos en deux catégories: avec animal et sans animal} \delM{Cet outil disponible gratuitement repose sur l'architecture YOLO version 5 basée sur des réseaux de neurones convolutionnels et a été entraîné par Microsoft à reconnaître différents taxons. MegaDetector n'identifie pas les éléments de l'image à l'espèce, mais catégorise plutôt les éléments comme étant des animaux ou non. Pour chaque élément identifié dans une image, MegaDetector lui attribue une probabilité à être dans la bonne catégorie (i.e., animal ou non). Nous avons lancé MegaDetector sur un serveur de calcul dédié du laboratoire du professeur Mazerolle à l'Université Laval.}

\delM{Afin de réduire le nombre de faux positifs, nous avons filtré les images pour lesquelles la probabilité était $\geq 75$\%. Nous avons visionné chacune de ces images afin de valider si la détection était avérée et le cas échéant, nous avons identifié les individus sur les images au taxon le plus bas possible (famille, genre ou espèce). Ces informations ont ensuite été utilisées pour évaluer l'occurrence de différents taxons à proximité des ponceaux. De plus, ces images annotées serviront à l’entraînement d’algorithmes d’intelligence artificielle pour la reconnaissance à l'espèce, notamment à l’aide du logiciel YOLO version 11.}


\subsubsection{Enregistrements acoustiques}

Nous avons installé 18 enregistreurs acoustiques automatisés à la fin avril 2025 aux milieux humides près des tronçons, présentant des habitats favorables à l’occupation \addM{par} des anoures (AudioMoths; \citealt{hill2019}). Ces appareils étaient placés à environ 1,5 m du sol et fixés à des arbres ou à des poteaux. Les appareils ont été programmés pour enregistrer les sons durant trois minutes, à cinq reprises dans la journée (20h00, 00h00, 7h00, 10h00, 14h00). Ces plages sont celles pendant lesquelles différentes espèces d’anoures sont actives, incluant la rainette faux-grillon (\emph{Pseudacris triseriata}) \citep{fayard2022}. Ces enregistreurs ont été retirés à la fin octobre 2025, à l'exception des enregistreurs sur T20, T48, T49\_1 et T57 qui ont été retirés à la fin juillet ou en août 2025 (Tableau \ref{tab:tab_AM25}).

Les plages d'enregistrement ont été analysées à l’aide de BirdNET, un outil de reconnaissance automatique utilisant l’intelligence artificielle \citep{kahl2021}. Cet outil permet d’identifier sur chaque enregistrement sonore les espèces pour lesquelles il a été entraîné. Les enregistreurs acoustiques nous ont permis de suivre les périodes de reproduction des anoures aux sites où l’accès est plus difficile pour le suivi visuel (longue distance à faire à pied, stationnement interdit en bordure d’autoroute), mais également de déterminer les espèces qui se reproduisent à proximité des tronçons à l'étude.

Bien que le modèle BirdNET soit déjà entraîné à reconnaître dix des espèces d'anoures présentes au Québec, deux espèces ne sont pas prises en charge par l’algorithme : la grenouille du Nord (\emph{Lithobates septentrionalis}) et la grenouille léopard (\emph{Lithobates pipiens}). Afin d’assurer une détection exhaustive des anoures, nous avons donc utilisé une version personnalisée de BirdNET que nous avons entraînée dans le laboratoire du professeur Mazerolle pour ces deux taxons. Cet entraînement reposait sur un corpus de 2 142 extraits acoustiques d’une durée de trois secondes chacun, préalablement sélectionnés et validés à la main. Les données provenaient de plateformes de science participative, notamment Xeno-canto, iNaturalist et Macaulay Library et extraites à l’aide du package R suwo \citep{arya-salas2025suwo}. Nous avons complété ce corpus d'entraînement avec des enregistrements issus de travaux antérieurs du laboratoire de biologie de la conservation quantitative du professeur Mazerolle. Une attention particulière a été portée à l’équilibrage du jeu d’entraînement afin de limiter les biais et le surapprentissage associés à certaines catégories de sons. Les 2 142 extraits étaient répartis en neuf classes distinctes (Tableau~\ref{tab:tab_clips}).


\begin{table}[H]
\centering
\caption{Composition du corpus d'entraînement du modèle de classification personnalisé BirdNET. Chaque clip a une durée de 3 secondes. Les clips ont été extraits à partir de fichiers audios issus des bases de données de sciences participatives (Xeno-canto, iNaturalist et Macaulay) et de la collection sonore du laboratoire du professeur Mazerolle.}
\label{tab:tab_clips}
\begin{tabular}{lrl}
  \hline
Classe & Nombre de clips (3s) & Description  \\ 
  \hline
ANTHRO & 124 & Bruits d'origine anthopique (moteurs et véhicules) \\
BIRDS & 247 &  Vocalisation d'oiseaux \\
CLAMITANS & 250 &  \emph{Lithobates clamitans} (Grenouille verte) \\
CRUCIFER & 250 &  \emph{Pseudacris crucifer} (Rainette crucifère) \\
PIPIENS & 242 &  \emph{Lithobates pipiens} (Grenouille léopard) \\
SEPTENTRIONALIS & 284 & \emph{Lithobates septentrionalis} (Grenouille du Nord) \\
SILENCE & 245 &  Clips sans signal biologique \\
VERSICOLOR & 250 & \emph{Dryophytes versicolor} (Rainette versicolore) \\
WEATHER & 250 &  Bruits météorologiques (vent et pluie) \\
     \hline
\end{tabular}
\end{table}

Nous avons évalué les performances du modèle personnalisé à partir d'un ensemble indépendant de 386 clips audio de 3 secondes, représentant environ 20 \% du corpus d'entraînement. \delM{Notons que cette proportion ne correspond à aucun seuil documenté.} \delM{Elle résulte simplement}\addM{Ceci correspondait à} un compromis entre \addM{le} nombre d'extraits suffisant pour obtenir des résultats concluants et \addM{la disponibilité} des données accoustiques. Cet ensemble comprenait des clips de grenouille du Nord et de grenouille léopard, ainsi que des clips négatifs (bruits environnementaux sans grenouille du Nord ou grenouille léopard) afin de tester la capacité du modèle à discriminer les deux espèces cibles du bruit de fond. Afin de garantir l’indépendance de l’évaluation, nous avons veillé à ce que ces clips soient issus de fichiers qui n’avaient pas servi à entraîner BirdNET en amont. \delM{Toutefois,}\addM{La validation sur un jeu de données délimité est un premier pas, mais ne reflète} pas pleinement le comportement réel du modèle en situation réelle. \addM{En effet, une meilleure évaluation devrait utiliser un plus grand volume de données}\delM{où ce dernier est exposé à un volume de données considérablement plus important}, à une plus grande variabilité acoustique et à des sons pour lesquels il n'a pas été entraîné.\delM{, augmentant ainsi le risque de faux positifs.} Il est donc attendu que notre modèle personnalisé produise encore un nombre élevé d'erreurs à ce stade. Ce taux devrait diminuer au fur et à mesure de l'entraînement, l'objectif étant d'obtenir un outil suffisamment fiable pour réduire considérablement le temps consacré à l'identification des espèces dans les fichiers acoustiques.

Malgré la présence de faux positifs\addM{, le modèle a démontré des performances satisfaisantes pour les deux espèces cibles, avec une exactitude (proportion de prédictions vraies sur l’ensemble des prédictions) supérieure à 0,70 pour les deux seuils testés, particulièrement pour la grenouille Nord} (Tableau~\ref{tab:tab_confusion_pipiens}; Tableau~\ref{tab:tab_metrics_pipiens_025}; Tableau~\ref{tab:tab_confusion_pipiens_075}; Tableau~\ref{tab:tab_metrics_pipiens_075}; Tableau~\ref{tab:tab_confusion_septentrionalis_025}; Tableau~\ref{tab:tab_metrics_septentrionalis_025}; Tableau~\ref{tab:tab_confusion_septentrionalis_075}; Tableau~\ref{tab:tab_metrics_septentrionalis_075}). Ces résultats indiquent que le modèle peut être utilisé comme outil de présélection des détections, réduisant substantiellement le temps de validation par rapport à l'écoute intégrale des enregistrements\addM{, à condition que les détections retenues fassent l'objet d'une validation manuelle}. Les faux positifs observés s'expliquent en partie par la similarité entre les vocalisations de \emph{L. septentrionalis} et celles d'autres espèces, notamment \emph{L. sylvaticus}. La constitution d'un corpus d'entraînement plus étoffé, incluant des enregistrements printaniers et estivaux de 2026, permettront d’améliorer les performances du modèle.

Avant l’analyse des fichiers audios, nous avons exclu les enregistrements affectés par des conditions météorologiques défavorables (vents forts et précipitations intenses) afin de préserver la qualité des données. Pour ce faire, nous avons utilisé notre modèle personnalisé préalablement entraîné pour la grenouille du Nord et la grenouille léopard, capable de détecter la présence de vent et de pluie à partir du signal acoustique (Tableau \ref{tab:tab_clips}). Nous avons utilisé ce modèle sur l'ensemble des enregistrements pour filtrer les enregistrements de mauvaise qualité. Ainsi, tout fichier contenant au moins un clip de 3 secondes classé dans la catégorie de mauvaise qualité (i.e., WEATHER, voir Tableau \ref{tab:tab_clips}) avec un score de confiance supérieur ou égal à 0,75 a été retiré des analyses subséquentes.

\subsubsection{Recherche par drone}

Comme en 2024, nous avons réalisé des suivis par drones au-dessus des zones humides de l'A10 et de la R104, afin de détecter les tortues qui thermorégulent dans les milieux humides \citep{aota2021, bogolin2021, bevan2018}. Nous avons utilisé le même drone décrit plus haut pour ces inventaires, en suivant le même protocole pour les observateurs visuels. \delM{Cette technologie permet d'éviter l'effarouchement des tortues, en respectant les recommandations de vol à une altitude minimale de 20 m au-dessus des milieux humides. Les survols en drone permettent aussi une couverture complète de la zone d’inventaire, et facilitent l’accès à des secteurs qui auraient été inaccessibles à pied.}

\delM{Transports Canada réglemente les zones où le vol de drone est autorisé. La zone de vol doit se situer à plus de 9 km de tout aéroport, héliport ou aérodrome. De même, les vols sont interdits au-dessus d’espaces réglementés, comme une base militaire, une prison ou une aire protégée comme celle du Bois de Brossard.} Les secteurs d’inventaire que nous avons sélectionnés \addM{n’étaient pas soumis à des interdictions de vols associées à la proximité d'espaces réglementés par Transports Canada}\delM{avaient pas de contrainte de vol}. \delM{Nous avons cherché des aires de décollage et d’atterrissage sécuritaires.} Ces suivis ont été réalisés à différentes périodes de la saison d'activité des tortues entre mai et octobre \delM{2024}\comM[inline]{C'était bien en 2025, non?}\addM{2025} \citep{boyer1965, obbard1981, lefevre1995}.

\delM{Nous avons suivi la méthodologie de recherche active proposée dans le protocole standardisé de détection et d'identification des tortues d'eau douce à l'aide de drones au Québec}\addM{Comme décrit dans le volet 1, nous avons suivi le protocole de recherche active proposé pour la détection et l'identification des tortues d'eau douce au Québec} \citep{mffp2021}. \delM{Nous avons rempli le "formulaire de prise de données" figurant dans l'annexe C de ce même protocole après chaque sortie. Les sorties se faisaient selon les conditions météorologiques suivantes : journée ensoleillée et couverture nuageuse pouvant aller jusqu’à 75 \%, température de l’air entre 10 et 25 \textdegree C, journée sans vent ou jusqu’à une légère brise (niveau 2 sur l’échelle de Beaufort). La caméra thermique du drone était utilisée pour connaître la température des tortues observées, leur substrat ainsi que l’eau du milieu où elles se trouvaient. Pendant une mission de vol, les milieux humides ont été survolés à basse vitesse à une altitude variant de 30 à 40 m. Lorsqu'une forme ressemblant à une tortue était repérée, nous avons utilisé le zoom du drone allant jusqu'à un grossissement de 56x pour confirmer qu'il s’agissait d'une tortue. En plus de la recherche active, lors d’une même visite nous avons lancé des vols automatisés qui filmaient l’intégralité du milieu humide en suivant toujours le même chemin à une altitude fixe de 40 m et un angle de caméra fixé à 90 \textdegree. Le but des vols automatisés était de constituer une banque de vidéos qui pourront servir à entraîner une intelligence artificielle (IA) pour identifier les tortues de manière automatique. Étant donné que la recherche active est plus longue que les vols automatisés,  nous voudrions utiliser l’IA sur les vols automatisés de tous les sites suivis.}



\subsubsection{Stations de plaques à reptiles}

Selon le même dispositif que celui décrit dans la méthode de la partie I, 22 stations de plaques à reptiles ont été installées et réparties parmi\delM{entre} les 30 tronçons entre le 5 et le 30 mai 2025. \delM{Nous n'avons pas pu attribuer une station de plaques à reptiles par tronçon, car parfois}\addM{Il a été impossible d'installer des plaques à reptiles dans certains tronçons où le milieu était défavorable aux couleuvres.} Sur les 22 \delM{stations installées}\addM{tronçons avec une station}\comM{Pas clair si 1 station = 1 tronçon. Il faut le dire clairement.}, trois visites ont été effectuées en 2025 : une au printemps, une durant l'été et une à la fin de l'été\comM[inline]{Ça prendrait des plages de dates pour quantifier \og printemps \fg{}, \og été \fg{} et \og fin d'été \fg{}.}. Les installations \delM{étaient}\addM{ont été} démantelées durant l'automne 2025.

